% ======================================================================
% APÉNDICE B: MANUAL DE INSTALACIÓN Y USO DE FATIGUESET-LIB
% ======================================================================

\chapter{Manual de Instalación y Uso de \code{fatigueset-lib}}
\label{ap:manual_libreria}

\section{Requisitos del Sistema e Instalación}
\label{sec:ap_instalacion}

La librería \code{fatigueset-lib} es compatible con sistemas operativos Linux (Ubuntu 20.04+, CentOS/RHEL 8+) y Windows 10/11 x86\_64. Requiere una versión de Python $\ge 3.10$ y soporte opcional para aceleración hardware mediante NVIDIA CUDA ($\ge 11.8$).

\subsection{Creación del Entorno Virtual e Instalación Local}
Se recomienda aislar las dependencias en un entorno virtual dedicado:

\begin{lstlisting}[style=bashstyle, caption={Instalación de \code{fatigueset-lib} en modo desarrollo editable.}]
# 1. Clonar el repositorio
git clone https://github.com/EmilioGullon/TFG_FatigueSet.git
cd TFG_FatigueSet/fatigueset-lib

# 2. Crear y activar entorno virtual
python -m venv .venv
source .venv/bin/activate  # En Windows: .venv\Scripts\activate

# 3. Instalar dependencias y paquete en modo editable
pip install --upgrade pip
pip install -e .

# 4. Verificar instalacion ejecutando los tests unitarios
pytest tests/ -v
\end{lstlisting}

\section{Guía Rápida de la API en Python}
\label{sec:ap_api_guia}

A continuación se ilustran ejemplos mínimos ejecutables de las principales funcionalidades del framework:

\subsection{Carga de Datos y Creación de DataLoader con GroupKFold}
\begin{lstlisting}[language=Python, caption={Carga de particiones estratificadas por participante.}]
from fatigueset.data import load_fatigueset_processed, get_group_kfold_loaders

# Cargar tensores preprocesados y metadatos de sujetos
X, y, groups = load_fatigueset_processed(data_dir="data/processed")

# Generar 5 particiones GroupKFold
for fold, (train_loader, val_loader) in enumerate(
    get_group_kfold_loaders(X, y, groups, n_splits=5, batch_size=32)
):
    print(f"--- Fold {fold+1} ---")
    print(f"Batches Train: {len(train_loader)} | Batches Val: {len(val_loader)}")
    break
\end{lstlisting}

\subsection{Instanciación y Entrenamiento de Modelos con la Factoría PyTorch}
\begin{lstlisting}[language=Python, caption={Entrenamiento de un modelo xLSTM con el motor desacoplado.}]
import torch
import torch.nn as nn
from fatigueset.models import get_model
from fatigueset.models.engine import train_model

device = torch.device("cuda" if torch.cuda.is_available() else "cpu")

# Instanciar el modelo xLSTM estabilizado
model = get_model(
    model_name="xlstm",
    input_size=23,
    hidden_size=160,
    num_layers=3,
    dropout=0.1,
    output_dim=2
).to(device)

criterion = nn.MSELoss()
optimizer = torch.optim.Adam(model.parameters(), lr=4e-3, weight_decay=1e-4)

# Entrenar con early stopping
history, best_model = train_model(
    model=model,
    train_loader=train_loader,
    val_loader=val_loader,
    loss_fn=criterion,
    optimizer=optimizer,
    epochs=50,
    patience=10,
    device=device
)
\end{lstlisting}
